\section*{一、预习报告}

\subsection*{1. 实验综述}

通过预习，我了解到液体表面层中的分子比内部稀疏，分子间作用表现为引力，使液面具有收缩的趋势，这就是表面张力（surface tension）。在液面上任取一条长为 $l$ 的分界线，两侧液面的相互拉力与分界线垂直、与液面相切，且满足

\[
f=\delta\, l,
\]

其中 $\delta$ 为表面张力系数，单位为 $\mathrm{N/m}$。它与液体的成分、纯度和温度有关，温度升高时 $\delta$ 减小。

本实验采用拉脱法（ring method）。实验中将把外径为 $D_1$、内径为 $D_2$ 的金属圆环挂在力敏传感器下并浸入水中，再缓慢降低液面，圆环内外两侧会带起一层液膜。液膜即将破裂时接触角 $\theta$ 很小，取 $\cos\theta\approx1$，并忽略浮力，则圆环所受拉力为

\[
F=mg+\pi\delta(D_1+D_2),\qquad \delta=\frac{F-mg}{\pi(D_1+D_2)}.
\]

力敏传感器由四个力敏电阻组成电桥，外力使电桥失去平衡，输出电压的变化与外力成正比，即 $k\Delta U=F-mg$，其中 $k$ 为传感器灵敏度（单位 $\mathrm{N/V}$）。因此

\[
\delta=\frac{k\,\Delta U}{\pi(D_1+D_2)}.
\]

由于圆环有一定厚度，液膜被拉起后先接近竖直，随后向内收缩，拉力在达到最大值后略有下降，最后在破裂瞬间骤降。所以电压差应取最大拉力对应的读数 $U_1$ 与液膜破裂后读数 $U_2$ 之差，即 $\Delta U=U_1-U_2$，而不是破裂瞬间前后的差值。

实验时，我将先打开变温界面张力分析仪充分预热（课件建议 15 分钟以上），并用水准泡将测试架调平。挂上砝码盘后将拉力置零，然后按 $0.500\,\mathrm{g}$ 逐个增加砝码至 $3.500\,\mathrm{g}$，再逐个减少至 $0$，分别记录电压。增减两次读数取平均，可以减小传感器滞后的影响；重力将按杭州 $g=9.79\,\mathrm{m/s^2}$ 换算，再作 $U$–$F$ 线性拟合，斜率的倒数即为 $k$。随后我将用游标卡尺在不同方向各测 6 次圆环外径 $D_1$ 和内径 $D_2$。

正式测量前，我将先用净水洗净圆环和加热器皿，再注入待测的水。挂上圆环并将拉力置零后，调节升降台使环底刚好接触液面，将位移置零，继续上升至位移约 $2\sim2.5\,\mathrm{mm}$，再缓慢下降，每隔 $0.05\,\mathrm{mm}$ 记录一次拉力电压，直至液膜被拉破，以观察拉力随位移的变化。同时我将用示波器（$200\,\mathrm{mV/DIV}$，时基 $2\,\mathrm{s/DIV}$ 或更大）记录拉力曲线，液膜破裂后按 STOP 暂停，用光标读取 $U_1$ 和 $U_2$。室温下计划重复测量 6 次；之后将用温控表使水温在室温基础上升高 $15\sim25\,{}^\circ\mathrm{C}$，待温度稳定后重复测量，并与标准值比较。

数据处理将按课件要求进行：由线性拟合得到 $k$ 及其标准差，并以 3 倍标准差估计 $k$ 的不确定度；$D_1$、$D_2$ 和 $\Delta U$ 将取多次测量的平均值并评定 A 类不确定度，游标卡尺的 B 类分量取其允许误差限 $\Delta=0.02\,\mathrm{mm}$，按均匀分布有 $u_B=\Delta/\sqrt3$。最后按传导法则合成

\[
\frac{u(\delta)}{\delta}=\sqrt{\left(\frac{u(k)}{k}\right)^{2}+\left(\frac{u(\Delta U)}{\Delta U}\right)^{2}+\frac{u^{2}(D_1)+u^{2}(D_2)}{(D_1+D_2)^{2}}}.
\]

电压读数的 B 类分量将以现场仪器给出的允许误差限为准，不能把显示分辨力直接当作允许误差限；本次预习不填入实测数值。

\subsection*{2. 实验目的}

\begin{enumerate}
\item 了解液体表面的性质，掌握拉脱法测量表面张力系数的原理和方法。
\item 学习力敏传感器的使用和砝码定标方法，掌握微小力的测量。
\item 测定水的表面张力系数，了解其随温度的变化，并学习不确定度的分析。

\end{enumerate}
\subsection*{3. 实验重点与难点}

\begin{enumerate}
\item 定标时要保持测试架水平，砝码轻放，待示数稳定后读数；增、减砝码各测一遍，检查传感器的线性和滞后。
\item 圆环和器皿必须洗净，圆环应保持水平；降低液面要缓慢，避免震动。读数时要正确识别最大拉力 $U_1$ 和破裂后的 $U_2$，不能过早读数，也不能取破裂瞬间前后的差值。
\item 加热后器皿温度较高，须防止烫伤；待水温稳定后再测量。实验结束先关加热电源，再关总电源，并用干布擦干仪器。

\end{enumerate}
